\section{Expérience}
  \resumeSubHeadingListStart
    \resumeSubheading
      {Développeur Full-Stack}{2024}
      {Boss Studios LLC}{À distance}
      \resumeItemListStart
        \resumeItem{Développé un jeu inspiré de Minecraft BedWars avec des mécaniques de jeu complètes}
        \resumeItem{Programmé le combat PvP, les systèmes de construction, les projectiles, les attaques au corps à corps et l'inventaire}
        \resumeItem{Intégré des systèmes de salons de jeu et des infrastructures multijoueurs}
      \resumeItemListEnd

    \resumeSubheading
      {Développeur Frontend}{2024}
      {Oakley Productions LLC}{À distance}
      \resumeItemListStart
        \resumeItem{Contribué à « Anime Defenders », un jeu de Tower Defense sur Roblox (+3,4 milliards de visites, pic de 100K connexions)}
        \resumeItem{Développé la logique de jeu pour plus de 10 unités/personnages et 30 capacités au total}
        \resumeItem{Conçu des effets visuels évolutifs en utilisant des systèmes de particules, des trajectoires de projectiles et des courbes de Bézier}
      \resumeItemListEnd

    \resumeSubheading
      {Développeur de Jeux Full-Stack}{2020 -- Présent}
      {Freelance/Contractuel}{À distance}
      \resumeItemListStart
        \resumeItem{Appliqué les mathématiques (matrices, vecteurs, cinématique inverse, calculs d'intersections) aux mécaniques de jeu}
        \resumeItem{Appliqué la physique (cinématique, mouvement de projectiles, ressorts, réflexion de lumière) pour des simulations réalistes}
        \resumeItem{Développé des systèmes réseau, animations personnalisées, effets visuels, UI/HUD, pathfinding IA et bases de données}
      \resumeItemListEnd
  \resumeSubHeadingListEnd

\section{Projets}
  \resumeSubHeadingListStart
    \resumeProjectHeading
      {\textbf{Projet d'Intégration du Cégep} $|$ \emph{Fusion360, CAO (Coll\`{e}ge Lionel-Groulx)}}{2026}
      \resumeItemListStart
        \resumeItem{Conçu une éolienne verticale de type H-Rotor à l'aide de logiciels de CAO}
        \resumeItem{Utilisé un moteur à courant continu (CC) à balais pour générer de la tension continue à partir de la rotation de l'éolienne}
        \resumeItem{Appliqué les principes de l'électromagnétisme, y compris la loi des mailles de Kirchhoff (KVL), la loi de Faraday, le courant induit et le couple d'encochage}
        \resumeItem{Appliqué les principes de l'ingénierie mécanique : aérodynamique, conception des profils d'aile, rapport d'engrenage et limite de Betz}
      \resumeItemListEnd

    \resumeProjectHeading
      {\textbf{Concours Déplace de l'air} $|$ \emph{Aérodynamique, Physique (Polytechnique Montréal)}}{2026}
      \resumeItemListStart
        \resumeItem{Dirigé une équipe de 3 personnes pour le concours de turbine éolienne}
        \resumeItem{Mené des recherches sur l'aérodynamique des pales, la conception des alternateurs et d'autres aspects physiques des éoliennes}
        \resumeItem{Collaboré avec les coéquipiers sur le remue-méninges, le prototypage et l'expérimentation}
        \resumeItem{Géré les contraintes de planification de projet, les compromis, les réunions et les tests du système}
      \resumeItemListEnd

    \resumeProjectHeading
      {\textbf{Simulation CFD} $|$ \emph{Simulation Physique, Moteur de Jeu, C\#}}{2026}
      \resumeItemListStart
        \resumeItem{Programmé une simulation de dynamique des fluides (CFD) sous les contraintes d'un moteur de jeu}
        \resumeItem{Implémenté un algorithme de pilotage inspiré des boids et des forces de pression pour simuler des trajectoires de particules}
        \resumeItem{Appliqué le principe de Bernoulli pour modéliser les relations de vitesse et de pression en temps réel}
      \resumeItemListEnd
  \resumeSubHeadingListEnd

\section{Activités Parascolaires}
  \resumeSubHeadingListStart
    \resumeSubheading
      {Mentor en Développement de Jeux \& Programmation}{2023 -- Présent}
      {Mentorat entre pairs / Tutorat}{À distance}
      \resumeItemListStart
        \resumeItem{Guidé et tutoré plusieurs pairs dans des cours de programmation informatique}
        \resumeItem{Mentoré un étudiant de fin de cycle pour son projet d'intégration Unity en 200.C1 (2025)}
        \resumeItem{Enseigné des concepts de développement de jeux : cinématique inverse FABRIK, pathfinding A*, mouvement de projectiles, logique de collision}
      \resumeItemListEnd

    \resumeSubheading
      {Bénévole aux Portes Ouvertes}{Coll\`{e}ge Lionel-Groulx}{Sainte-Th\'{e}r\`{e}se, QC}
      {2025}
      \resumeItemListStart
        \resumeItem{Accueilli et guidé les visiteurs dans le programme de sciences pures et appliquées sur le campus}
        \resumeItem{Répondu aux questions sur les cours, la charge de travail académique et les attentes}
      \resumeItemListEnd
  \resumeSubHeadingListEnd

\section{Compétences Techniques}
  \begin{itemize}[leftmargin=0.15in, label={}]
    \small{\item{
      \textbf{Langages}{: C\#, Python, JavaScript, TypeScript, Lua, SQL (Postgres + SQLite), HTML/CSS} \\
      \textbf{Frameworks}{: React, FastAPI, Next.js, Playwright, Vitest, React Native, Tailwind} \\
      \textbf{Outils de Développement}{: Git, GitHub, Vercel, Railway, Expo, Supabase, Redis, VS Code, Visual Studio, Vite, Fusion360, Unity, Blender, Roblox Studio} \\
      \textbf{Outils d'IA}{: Claude Code, Antigravity, Hermes, Gemini CLI, OpenCode, Goose} \\
      \textbf{Bibliothèques}{: PyTrends, NumPy, Matplotlib, Mammoth, PDF Parse}
    }}
  \end{itemize}
